\documentclass[main.tex]{subfiles}
\begin{document}\newpage
\setdoublesep{0.35700 em}  % 'Bond Spacing'
\setatomsep{1.78500 em}    % 'Fixed Length'
\setbondoffset{0.18265 em} % 'Margin Width'
\newcommand{\bondwidth}{0.06642 em} % 'Line Width'
\setbondstyle{line width = \bondwidth}

\begin{fullwidth}





%%%%%%%%%%%%HEADING
\begin{multicols}{2}
\begin{tcolorbox}[enhanced jigsaw,breakable,size=title,
colback=mybrown!05,colframe=black,fonttitle=\bfseries,
title=STUDENT INFO,pad at break=1mm, break at=15cm/0pt ]
\vspace{0.2cm}
\noindent Name: \rule{5cm}{0.4pt}Date:\rule{1cm}{0.4pt}\\
Pre-lab Done: \tikzcheckmark[scale=2,black]{no mark}\quad
\end{tcolorbox}
\end{multicols}
\hfill
\vspace{0.2cm}
\begin{center}
{\large \bfseries 
Pre-lab Questions 
\par
\Huge
Solubility
\\[5pt] \par}
\vspace{0.2cm}
\end{center}
\par
\noindent
\uline{  \hfill \normalsize \hfill       }
%%%%%%%%%%%%HEADING

\begin{enumerate}
% PELAB 1
\item Name and predict the solubility (\textit{soluble} or \textit{insoluble}) of the following substances.\vspace{0.8cm} 
\tabulinestyle{1 pt black }
\begin{center}\begin{tabu} to 150mm {XXX}
 \tabucline-
Compound & Name &Solubility\vspace{0.8cm}     \\
 \tabucline-
\ce{NaNO3} \vspace{0.8cm} &&    \\ 
\tabucline-
\ce{NH4Br} \vspace{0.8cm} & &   \\ 
\tabucline-
\ce{BaBr2} \vspace{0.8cm} &&    \\ 
\tabucline-
\ce{Ni(OH)2} \vspace{0.8cm} & &    \\ 
\tabucline-
\ce{BaCO3} \vspace{0.8cm} &&    \\ 
\tabucline-
\end{tabu}\end{center}
\vspace{0.8cm}

\item Are salts more or less soluble in liquid at high or low temperature?
\vspace{3cm}

\item Think about a soda can. Are gases more or less soluble in liquid at high or low temperature?
\vspace{3cm}



 

\end{enumerate}


\clearpage\mbox{}\clearpage



%%%%%%%%%%%%HEADING
\begin{multicols}{2}
\begin{tcolorbox}[enhanced jigsaw,breakable,size=title,
colback=mybrown!05,colframe=black,fonttitle=\bfseries,
title=STUDENT INFO,pad at break=1mm, break at=15cm/0pt ]
\vspace{0.2cm}
\noindent Name: \rule{5cm}{0.4pt}Date:\rule{1cm}{0.4pt}\\
Pre-lab Done: \tikzcheckmark[scale=2,black]{no mark}\quad
\end{tcolorbox}
\end{multicols}
\hfill
\vspace{0.2cm}
\begin{center}
{\large \bfseries 
Experiment
\par
\Huge
Solubility\\[5pt] \par}
\vspace{0.2cm}
\end{center}
\par
\noindent
\uline{  \hfill \normalsize \hfill       }
%%%%%%%%%%%%HEADING

\vspace{0.2cm}{\large \bfseries 1. Solubility rules}
The goal of this mini experiment is to test the solubility rules and predict the soluble character of a chemical.
\begin{steps}[resume]
\newstep[] Find a spot plate. Arrange in the following order the set of reactants Type A:  \ce{NaCl(aq)}, \ce{Na2SO4(aq)} and \ce{Na3PO4(aq)}. Arrange in the following order the set of reactants Type B:  \ce{Ca(NO3)2(aq)} and \ce{AgNO3(aq)}.
\newstep[] Make mixtures of each pair of compounds listed in the results table by adding 2-3 drops of each solution in the same spot of the spot plate. When the resulting mixture is cloudy that means a precipitate has formed. 
\newstep[] Write down the result on the Results table as soluble (S) or insoluble (I).

 \end{steps} 
     
     \begin{center}\begin{tabular}{ |p{2cm}|p{2cm}|p{2cm}|p{2cm}|  }
\cline{2-4}
\multicolumn{1}{r|}{} & \ce{NaCl} \hspace{0.3cm}& \ce{Na3PO4(aq)}\hspace{0.3cm}& \ce{Na2SO4}\hspace{0.3cm} \\ [7pt]
\hline
\vspace{0.3cm}\ce{Ca(NO3)2(aq) }&    & &   \\
\hline
\vspace{0.3cm}\ce{AgNO3(aq) } &    & &   \\
\hline
\end{tabular}\end{center}
\begin{center}(write \textit{S} for soluble product and \textit{I} for insoluble product.)\end{center}


\begin{figure*}[h] % FUL FIGURE
        \includegraphics[width=\linewidth,scale=0.08]{./chapter8/spot.png}
\caption{A spot plate used for this miniexperiment.}
\end{figure*}
\vspace{1cm}


\end{fullwidth}



\newpage 
\begin{fullwidth}
\vspace{0.2cm}{\large \bfseries 2. Change of solubility with temperature}
The goal of this mini experiment is to investigate the change of solubility of \ce{KNO3} with temperature. You will do so by adding different amounts of solute and measuring the temperature at what the solute dissolved. You will work in different teams and each team will share their results with the rest of the class.
\begin{steps}
    \newstep[]  You will be assigned an amount of solute between 3 and 7 grams. Weight the solid and write down exactly how much solute did you use. If for example you are assigned 3g you can weight 3.1g.
        \newstep[]  Measure 5mL of water with a measuring cylinder. Add the liquid and the solid to a test tube that should be  camped to a stand inside a water bath. Use a thermometer and start warming up the water bath.
          \newstep[]  Heat the solution until the solid dissolves completely. At that point, stop the heating and take the tube out of the bath keeping the thermometer inside the solution. When crystals start to appear, write down the temperature of the solution. You can calculate solubility in $g\cdot mL^{-1}$ using the formula:
          \[\text{Solubility}=\frac{\text{mass of \ce{KNO3}}}{\text{5 mL of \ce{H2O}}}\times 100\]
           \newstep[] When you have all results from the class, plot solubility (vertical axis) vs. temperature (horizontal axis).
       \end{steps}
\vspace{2cm}




   \begin{center}\begin{tabular}{ |>{\centering\arraybackslash}p{3cm}|>{\centering\arraybackslash}p{6cm}|>{\centering\arraybackslash}p{3cm}|  }
\hline
 Mass of \ce{KNO3}  (g) & Temperature when crystals appear ($^\circ$C)& Solubility (g/ml)  \\ [7pt]
\hline
3\vspace{0.3cm}      & &   \\
\hline
3.5 \vspace{0.3cm}   & &   \\
\hline
 4  \vspace{0.3cm}   & &   \\
\hline
 4.5 \vspace{0.3cm}   & &   \\
\hline
 5  \vspace{0.3cm}   & &   \\
\hline
 5.5  \vspace{0.3cm}   & &   \\
\hline
 6 \vspace{0.3cm}   & &   \\
\hline
 6.5  \vspace{0.3cm}   & &   \\
\hline
 7  \vspace{0.3cm}   & &   \\
\hline
 \end{tabular}\end{center}




\end{fullwidth}


\newpage
\begin{fullwidth}
\begin{center}
 \begin{tikzpicture}
 \draw[step=1mm,help lines] (0,0) grid (120mm,120mm);
 \draw[step=10mm] (0,0) grid (120mm,120mm);
\end{tikzpicture}
\end{center}


\vspace{0.6cm}{\large \bfseries PostLab questions }
\begin{enumerate}
\item[2 a)] According to your graph, estimate solubility at 55 $^\circ$C.
\vspace{3.7cm}
\item[2 b)] Solubility for a given chemical is 0.1 $g\cdot mL^{-1}$ at 30 $^\circ$C. How many grams of solute will dissolve in 25mL of water at that temperature.
\vspace{3.7cm}
\item[2 c)] Solubility of table salt is 0.4 $g\cdot mL^{-1}$ at 25 $^\circ$C. Will 50 grams of table salt dissolve in 100mL of water at that temperature.
\vspace{3.7cm}
\end{enumerate}


\end{fullwidth}




\end{document}